\documentclass[a4paper,10pt]{article}
\usepackage[T1]{fontenc}
\usepackage[utf8]{inputenc}
\usepackage[dutch]{babel}
\usepackage[margin=20mm]{geometry}
\usepackage{array,longtable,booktabs}
\newcolumntype{P}[1]{>{\raggedright\arraybackslash}p{#1}}
\setlength{\parindent}{0pt}
\renewcommand{\arraystretch}{1.25}
\begin{document}
\section*{Bijlage A -- Pinverbindingen}
De nummers bij J1 en J2 zijn connectorpinnen uit PCB\_schema\_V1, niet de pinnummers van de losse IC's. Controleer bij een andere module de opschriften en pinvolgorde. De Pico-kolom noemt de fysieke headerpin en waar nodig het GPIO-nummer.

{\small
\begin{longtable}{@{}P{4.7cm}P{5.1cm}P{6cm}@{}}
\toprule
\textbf{Van / aansluiting} & \textbf{Naar / aansluiting} & \textbf{Functie} \\
\midrule
\endfirsthead
\toprule
\textbf{Van / aansluiting} & \textbf{Naar / aansluiting} & \textbf{Functie} \\
\midrule
\endhead
Pico pin 36 (3V3) & J1-1 (VDD), J2-8 (VCC) & Voeding van beide modules: 3,3 V. \\
Pico GND (bijv. pin 38) & J1-2, J1-5, J2-7, Q1-3 en PV$-$ & Gemeenschappelijke massa; J1-5 is A0. \\
Pico pin 26 (GPIO20) & J1-4 (SDA), J2-6 (SDA) & I$^2$C-data; GPIO20 is niet fysieke pin 20. \\
Pico pin 27 (GPIO21) & J1-3 (SCL), J2-5 (SCL) & I$^2$C-klok. \\
J1-6 (DAC OUT) & Q1 pin 1 (gate) & Analoge stuurspanning voor de belasting. \\
PV+ (zonnecel) & J2-1 (IN+) en J2-3 (VBS) & Positieve stroomingang en busspanningsmeting. \\
J2-2 (IN$-$) & Q1 pin 2 (drain) & Stroompad na de shunt op de module. \\
Q1 pin 3 (source) & GND / PV$-$ & Retour van de belastingsstroom. \\
J2-4 (ALE / ALERT) & Niet aangesloten & Alarmuitgang niet gebruikt. \\
R1 (4,7 k$\Omega$) & Tussen 3V3 en SDA & Pull-up voor de datalijn. \\
R2 (4,7 k$\Omega$) & Tussen 3V3 en SCL & Pull-up voor de kloklijn. \\
C1 (100 nF) & Tussen J1-1 (VDD) en GND & Ontkoppeling van de DAC-voeding. \\
C2 (100 nF) & Tussen J2-8 (VCC) en GND & Ontkoppeling van de INA226-voeding. \\
Computer USB & USB-aansluiting Pico & Voeding en uitwisseling van opdrachten en meetwaarden. \\
\bottomrule
\end{longtable}
}
VBS is de moduleaansluiting voor VBUS. De tabel volgt de verbinding met IN+ aan PV+. Controleer deze verbinding op de gebruikte module; alleen dan geldt zonnecelspanning = gemeten busspanning. Bij Q1 gelden gate = 1, drain = 2 en source = 3; controleer de montagerichting met de datasheet.
\clearpage
\section*{Bijlage B -- Materialenlijst}
Benodigdheden per meetopstelling. Bij een reeds bestukte PCB zijn de Pico, modules en losse printonderdelen inbegrepen; deze worden niet dubbel uitgedeeld. Gebruik de door de begeleider vrijgegeven stroom- en vermogensgrenzen.

{\small
\begin{longtable}{@{}P{1.5cm}P{6cm}P{8.3cm}@{}}
\toprule
\textbf{Aantal} & \textbf{Materiaal} & \textbf{Toepassing} \\
\midrule
\endfirsthead
\toprule
\textbf{Aantal} & \textbf{Materiaal} & \textbf{Toepassing} \\
\midrule
\endhead
1 & Zonnecel of zonnecelmodule & Met aansluitdraden naar PV+ en PV$-$. \\
1 & Lamp met voeding en houder & Voor een vaste, zo homogeen mogelijke belichting. \\
1 & Zonnecel-PCB volgens PCB\_schema\_V1 & Bestukte print met de hieronder genoemde onderdelen. \\
1 & Raspberry Pi Pico & Met passende headers en meetfirmware. \\
1 & MCP4725 DAC-module & 12-bit DAC op connector J1. \\
1 & INA226-meetmodule met shunt & Op J2; noteer de werkelijke shuntwaarde voor de kalibratie. \\
1 & IRLZ44N MOSFET (Q1) & Met koeling passend bij de vrijgegeven dissipatie. \\
2 & Weerstand 4,7 k$\Omega$ & R1 en R2: I$^2$C-pull-ups. \\
2 & Condensator 100 nF & C1 en C2: voedingsontkoppeling. \\
1 set & Headers, aansluitklemmen en bedrading & Voor de modules en de zonnecelaansluiting. \\
1 & Computer met meet- en analysesoftware & Voor instellen, opslaan en verwerken van de meetreeks. \\
1 & USB-datakabel voor de Pico & Een kabel die ook data overdraagt. \\
1 & Multimeter & Voor controle van verbindingen en spanning. \\
1 set & Houders en liniaal of meetlint & Voor vaste afstand en stand van lamp en zonnecel. \\
Optioneel & Gekalibreerde instralingsmeter & Nodig voor rendement uit instraling in W/m$^2$ en celoppervlak. \\
\bottomrule
\end{longtable}
}
\end{document}
